% =========================================================================
% APPENDIX B — STUDY PROTOCOLS
% =========================================================================

\chapter{Study Protocols}
\label{app:protocols}

This appendix documents the instruments of the empirical evaluation. The study is delivered as a self-guided, self-administered online questionnaire, so the protocol is the questionnaire itself: its structure, its four variants, its comprehension items, the consent and information text that precede it, and the coding grid against which the responses are analysed. The full text of the four questionnaire variants is reproduced in \S\ref{app:b-fulltext}; the analytical key that governs how their open responses are read is set out in \S\ref{sec:app-coding-grid}.

% -------------------------------------------------------------------------
\section{Questionnaire Variants}
\label{app:b-variants}

The instrument exists in four variants, crossing two languages with two conditions:

\begin{itemize}
    \item \textbf{Language:} Italian and English, so that participants respond in their first language. Open responses are analysed in the language in which they are written (Appendix~\ref{app:limits-evaluate}).
    \item \textbf{Condition:} a \emph{control} version showing only the conventional static representation (C1), and an \emph{experimental} version showing the dynamic and experiential representation (C2 and C3).
\end{itemize}

\noindent The four variants (IT-control, IT-experimental, EN-control, EN-experimental) are identical in structure, background questions, and comprehension items. They differ at exactly four points: the \textsc{what you'll do} section, the three scene-opening exploration blocks, the Scene~III controls paragraph, and the group-reveal label. Every other section is held constant so that the information conveyed to the two groups is identical and only the representational form varies, preserving informational equivalence \parencite{tversky2002} (Appendix~\ref{app:limits-evaluate}).

% -------------------------------------------------------------------------
\section{Structure of the Questionnaire}
\label{app:b-structure}

Each participant follows the same journey, in four stages.

\paragraph{1. Consent and information.}
The questionnaire opens with the plain-language information sheet and the consent statement (\S\ref{app:b-consent}). Participation proceeds only after consent is given.

\paragraph{2. Background questions.}
A short set of questions characterises the participant and allows prior knowledge to be accounted for in analysis:
\begin{itemize}
    \item age band;
    \item highest level of education;
    \item field of study or work;
    \item self-rated familiarity with climate science, with the AMOC specifically, and with reading data visualisations.
\end{itemize}
These items collect no directly identifying information; responses are anonymised (Appendix~\ref{app:ethics-conduct}).

\paragraph{3. Guided exposure with per-scene comprehension.}
The participant is taken through the three scenes of their assigned condition in sequence. Each scene opens with a neutral introduction carrying its driving question, links to the assigned artefact, and on return poses its comprehension questions before the next scene begins, so that understanding is captured for each scene as it forms rather than reconstructed at the end. The comprehension items are described in \S\ref{app:b-comprehension}.

\paragraph{4. Final reflection and debrief.}
After the three scenes, a short reflection asks the participant to step back and consider the system as a whole. Only then does the debrief reveal the two-condition structure and disclose which condition the participant was in. At this point every participant --- including the control group --- is given access to the complete set of materials, the other condition included. This closing stage is unassessed and does not enter the between-group comparison; it serves as disclosure and as a return to the participant, and is placed outside the assessed journey by design (Appendix~\ref{app:limits-evaluate}).

% -------------------------------------------------------------------------
\section{Comprehension Items}
\label{app:b-comprehension}

The comprehension items are the study's primary evidence. Each scene poses three open, symmetric description questions --- what the participant saw, how the system behaves over depth or time, and what surprised or struck them --- followed by a single closed self-rating of confidence. The open questions target the structural property each scene is designed to convey, so that they measure understanding of what the artefact actually encodes rather than general climate knowledge.

The open format is deliberate: it reveals a participant's mental model in their own words rather than testing their ability to recognise a supplied answer. The closed confidence item is asked only after the free description it refers to, never before, so that rating one's certainty cannot reshape the account (\S\ref{sec:app-coding-grid}). The open responses are the primary material for the analysis; the confidence ratings are read only in cross-tabulation against whether the structure was grasped. No structural reading is ever stated to the participant in either condition: the instrument teaches how to read each scene without supplying what it shows.

% -------------------------------------------------------------------------
\section{Full Text of the Questionnaire Variants}
\label{app:b-fulltext}

The four variants are reproduced below in the form administered, aligned so that the four points of difference identified in \S\ref{app:b-variants} are visible against the shared common text. The live forms and their exported responses are held in the study repository documented in Appendix~\ref{app:code}.

% =========================================================================
% >>> ENRICO: PASTE THE VERBATIM TEXT OF THE FOUR QUESTIONNAIRE VARIANTS HERE <<<
% IT-control / IT-experimental / EN-control / EN-experimental.
% Suggested form: \subsection*{Shared text} then \subsection*{Points of difference}.
% Chapter 5 and Chapter 8 both promise the reader the full instrument at this point.
% =========================================================================

% -------------------------------------------------------------------------
\section{Information Sheet and Consent}
\label{app:b-consent}

The information sheet states, in plain language: what the study involves and how long it takes; that participation is voluntary and may be stopped at any point without consequence; what data are collected (background questions and questionnaire responses) and that no directly identifying information is gathered; how responses are anonymised, stored securely under the General Data Protection Regulation, retained only as long as needed and used only for this research; and contact details for the researcher and the supervisor. Consent is given explicitly before the questionnaire begins.

% =========================================================================
% >>> ENRICO: PASTE THE FULL INFORMATION SHEET AND CONSENT TEXT HERE (IT + EN) <<<
% =========================================================================

The ethical commitments these documents implement are set out in Appendix~\ref{app:ethics-conduct}.

% =========================================================================
\section{The Coding Grid}
\label{sec:app-coding-grid}

The open responses were analysed against the grid set out below, fixed before the responses were read. Its categories are of two kinds. The \emph{structural targets} record whether a participant grasped the claim each scene was built to deliver, and are summarised in Table~\ref{tab:coding-targets}. The \emph{diagnostic codes} record specific misreadings and cross-cutting properties of the response, and are defined in the codebook that follows. Nothing in the grid is shown to participants; it governs the analysis, not the instrument.

The reading is deductive at the level of these categories and interpretive within them: the grid fixes what is looked for, while judgement decides whether a given description, in its own words, meets a category's definition. Every code is applied to the three symmetric description questions each scene poses --- what the participant saw, how the system behaves, and what struck them --- and never to a closed item.

% -------------------------------------------------------------------------
\subsection{Structural Targets}
\label{subsec:app-targets}

Each scene carries one structural claim, named by its driving question and defined in \S\ref{sec:three-scenes}. A description is coded \emph{grasped} only when it names that structure in the participant's own words; the presence of movement, colour or general interest is not sufficient. Intermediate credit is recorded as \emph{partial} where one necessary component is present and another absent, as specified per scene.

\begin{table}[htbp]
\centering
\caption[Structural coding targets by scene]{The structural target for each scene, with the conditions under which a free description is coded \emph{grasped}, \emph{partial} or \emph{not grasped}. The targets correspond one-to-one to the driving questions of \S\ref{sec:three-scenes}.}
\label{tab:coding-targets}
\small
\renewcommand{\arraystretch}{1.35}
\begin{tabularx}{\textwidth}{@{}p{2.5cm} X X@{}}
\toprule
\textbf{Scene} &
\textbf{Structural claim (target)} &
\textbf{Coded \emph{grasped} when\ldots} \\
\midrule

\makecell[l]{Scene I\\The Column} &
Water moves in \emph{opposing} directions at different depths: broadly northward near the surface, southward at depth. &
the description names (b) two opposing directions \emph{and} (c) ties direction to depth. \emph{Partial}: direction mentioned (a) but not both opposing and depth-linked. \\

\makecell[l]{Scene II\\The Engine Room} &
The system \emph{revisits} states it has already occupied; motion is recurrent, not a monotonic drift toward the new. &
the description conveys return, cycling or passing back through prior states. \emph{Partial}: motion or change noted without the idea of return. \\

\makecell[l]{Scene III\\The Threshold} &
The \emph{agreement} between the three models \emph{changes} over time --- more disagreement earlier, convergence later. &
the description names a change in agreement across time. \emph{Partial}: agreement or disagreement noted as constant, without change over time. \\

\bottomrule
\end{tabularx}
\end{table}

% -------------------------------------------------------------------------
\subsection{Cross-cutting Axes}
\label{subsec:app-axes}

Four axes are recorded for every response, independent of the structural target, and are the basis for the descriptive counts reported in Chapter~\ref{ch:results}.

\begin{description}
\item[\textnormal{\textsc{Structure-grasped}}] The verdict from Table~\ref{tab:coding-targets}: grasped, partial or not grasped. The primary outcome, counted per condition.

\item[\textnormal{\textsc{Confidence}}] The participant's self-rating, recorded from the closed item that follows each free description. Read only in cross-tabulation, never as an outcome in itself.

\item[\textnormal{\textsc{Confidence\,$\times$\,correctness}}] The crossing of the two above. The cell of interest is high confidence with structure \emph{not} grasped --- the confidently wrong reading that \textcite{fischer2020} document, and that expectation~E3.2 concerns.

\item[\textnormal{\textsc{Uncertainty-as-quality}}] Whether the participant expresses uncertainty \emph{productively} --- naming what they are unsure of, or what a display does not tell them --- as distinct from either false confidence or blank non-response. Productive uncertainty is credited, not penalised; it is the positive face of the same axis whose negative face is the confidently wrong.
\end{description}

% -------------------------------------------------------------------------
\subsection{Diagnostic Codes}
\label{subsec:app-diagnostic}

Beyond the structural verdict, the grid records one misreading, identified in advance as likely and carrying interpretive weight. It is defined by what licenses the code and what excludes it, so that its application is repeatable.

\paragraph{S1-MISC-INT --- Streamfunction read as a velocity profile.}
\emph{Applies to: Scene~I.} The participant reads the \emph{height} of the $\Psi(z)$ curve as the speed or intensity of water movement at that depth, rather than reading its \emph{slope} as the local transport. Equivalently, $\Psi$ is treated as a direct velocity profile rather than as the cumulative integral of transport.

\emph{Positive indicators} (any one suffices): speed described as decreasing or increasing in step with the curve's height; ``no movement'' or stillness located at the $\Psi$ zero-crossing rather than at the slope-zero depths; the curve's peak read as maximum movement rather than as a flow reversal; the whole curve narrated as one continuous velocity story from surface to seabed.

\emph{Negative indicators} (do \emph{not} code): direction reversal correctly tied to where the curve changes slope sign; accumulated transport explicitly distinguished from local velocity; productive uncertainty about what the curve's height represents --- coded instead under \textsc{uncertainty-as-quality}.

\emph{Note.} This code is independent of the Scene~I structural target. A participant may name the opposing \emph{directions} correctly --- and so be coded \emph{grasped} --- while still misreading the curve's height as \emph{speed}, and thus also carry S1-MISC-INT. The two record different halves of the same reading and must not be collapsed: recording both is what reveals a participant who holds the direction of the flow but not the cumulative nature of the instrument. Crossed with condition, the code's frequency is the evidence for whether the experiential form reduces the misconception; crossed with confidence, for whether it is held with false certainty.

% -------------------------------------------------------------------------
\subsection{Coding Procedure}
\label{subsec:app-procedure}

The grid above was applied in two passes. In the first, a large language model was given the
grid, its category definitions and one response at a time, and returned a proposed label. In
the second, the researcher read every response against the grid and adjudicated the proposed
label, correcting it wherever the two readings differed. Every label in the reported dataset
is therefore the researcher's; the automated pass served to order the work, not to decide it.

Two consequences follow and are recorded rather than mitigated. There is no independent
second coder, so no inter-coder agreement coefficient is reported --- the measure does not
apply to a single-coder design. And a reviewer who sees a proposed label before forming their
own may anchor on it; the fixed criteria of the grid limit how far that can go, but do not
remove it. Both are listed among the limitations of the study in \S\ref{sec:eval-limitations}.

% -------------------------------------------------------------------------
\subsection{Emergent Codes}
\label{subsec:app-emergent}

Alongside the pre-specified grid, the coding recorded patterns that were not part of it.
These are post-hoc: they emerged during reading, they were not predicted, and they are
reported separately from the structural targets throughout Chapter~\ref{ch:results}. Three
recurred often enough to be reported by condition; the remainder occurred once each and are
listed for completeness without interpretation, since a single occurrence cannot be
distinguished from idiosyncrasy at this sample size.

\begin{description}
\item[\textnormal{\textsc{Interaction-aided}}] The response attributes its understanding to
acting on the display --- pausing, rewinding, isolating an element. Twelve occurrences, all
in the experimental condition.

\item[\textnormal{\textsc{Static-limit-named}}] The response names the absence of a temporal
ordering as the reason it cannot tell whether the system revisits earlier states. Four
occurrences, all in the control condition.

\item[\textnormal{\textsc{Anomaly-confusion}}] The response registers the usual/unusual
encoding but cannot say what it signifies. Six occurrences, one in the control condition and
five in the experimental.
\end{description}

% =========================================================================
% >>> ENRICO: LIST THE SINGLE-OCCURRENCE EMERGENT CODES HERE <<<
% Chapter 9 (footnote in \S\ref{sec:res-emergent}) promises the reader the full list.
% =========================================================================
